\documentclass[11pt,a4paper]{article}
\usepackage[margin=2cm]{geometry}
\usepackage{fontspec}
\usepackage{polyglossia}
\setdefaultlanguage{mongolian}
\setmainfont{cmunrm.otf}[BoldFont=cmunbx.otf,ItalicFont=cmunti.otf,BoldItalicFont=cmunbi.otf]
\usepackage{amsmath,amssymb}
\usepackage{array}

\setlength{\parindent}{0pt}
\setlength{\parskip}{3pt}

\newcommand{\bod}[1]{\par\medskip\noindent\textbf{#1}\quad}
\newcommand{\ol}[1]{\overline{#1}}

\begin{document}

\begin{center}
  {\Large Магадлалын онол. Гэрийн даалгавар 2}\\[4pt]
  Нэр: Дугаржанцан Билгүүн \qquad Код: 23B1NUM1470
\end{center}

\section*{Бодсон бодлогын хүснэгт}

I. О.Цэрэнбат, Магадлалын онол (2023), 2-р бүлэг

\smallskip
\setlength{\tabcolsep}{4pt}
\begin{tabular}{|l|*{14}{c|}}
\hline
Бодлого & 2.1 & 2.2 & 2.3 & 2.4 & 2.5 & 2.6 & 2.7 & 2.8 & 2.9 & 2.10 & 2.11 & 2.12 & 2.13 & 2.14 \\ \hline
 & + & + & + & + & + & + & + & + & + & + & + & + & + & + \\ \hline
Бодлого & 2.15 & 2.16 & 2.17 & 2.18 & 2.19 & 2.20 & 2.21 & 2.22 & 2.23 & 2.24 & 2.25 & 2.26 & 2.27 & 2.28 \\ \hline
 & + & + & + & + & + & + & + & + & + & + & + & + & + & + \\ \hline
\end{tabular}

\medskip
II. Bertsekas, Tsitsiklis, Introduction to Probability, 2-р бүлэг (х.~119--138)

\smallskip
\begin{tabular}{|l|*{11}{c|}}
\hline
Бодлого & 2 & 3 & 14 & 18 & 22 & 25 & 32 & 35 & 39 & 44 & 45 \\ \hline
 & + & + & + & + & + & + & + & + & + & + & + \\ \hline
\end{tabular}

%======================================================================
\section*{I. О.Цэрэнбат, Магадлалын онол, 2-р бүлэг}

\bod{2.1.}
Сүлдийг С, тоог Т гэж тэмдэглэвэл $\Omega=\{\text{СС}, \text{СТ}, \text{ТС}, \text{ТТ}\}$, эгэл үзэгдэл бүрийн магадлал $1/4$.
\[
A=\{\text{СС},\text{СТ}\},\quad D=\{\text{СС},\text{СТ},\text{ТС}\},\quad E=\{\text{СТ},\text{ТС},\text{ТТ}\},\quad F=\{\text{СС},\text{ТС}\},
\]
\[
P(A)=P(F)=\tfrac12,\qquad P(D)=P(E)=\tfrac34 .
\]
а) $P(A\cap E)=P(\{\text{СТ}\})=\frac14\ne P(A)P(E)=\frac38$ тул хамааралтай.

б) $P(A\cap F)=P(\{\text{СС}\})=\frac14=P(A)P(F)$ тул үл хамаарах.

в) $P(D\cap E)=P(\{\text{СТ},\text{ТС}\})=\frac12\ne\frac{9}{16}$ тул хамааралтай.

г) $D\cap F=F$, $P(D\cap F)=\frac12\ne\frac38$ тул хамааралтай.

д) $P(E/A)=\dfrac{1/4}{1/2}=\dfrac12,\quad P(A/F)=\dfrac{1/4}{1/2}=\dfrac12,\quad P(D/E)=\dfrac{1/2}{3/4}=\dfrac23,\quad P(D/F)=\dfrac{P(F)}{P(F)}=1.$

\bod{2.2.}
$P(A)=\frac{4}{52}=\frac{1}{13}$, $P(B)=\frac12$, $P(C)=\frac1{52}$, $P(D)=\frac1{13}$.

а) $A\cap B$ нь 2 улаан тамга: $P(A\cap B)=\frac{2}{52}=\frac1{13}\cdot\frac12$ тул үл хамаарах.

б) $C\subset A$ тул $P(A\cap C)=\frac1{52}\ne\frac1{13}\cdot\frac1{52}$, хамааралтай.

в) Бундан улаан өнгөтэй тул $C\subset B$, $P(B\cap C)=\frac1{52}\ne\frac1{104}$, хамааралтай.

г) $B\cap D$ нь 2 улаан 10-т: $P(B\cap D)=\frac2{52}=\frac12\cdot\frac1{13}$ тул үл хамаарах.

д) $C\cap D=\varnothing$, $P(C\cap D)=0\ne P(C)P(D)$ тул хамааралтай.

\bod{2.3.}
$n=a+b+c$ гэе. Ядаж хоёр нь ижил өнгөтэй байх үзэгдлийн эсрэг үзэгдэл нь гурвуулаа өөр өнгөтэй байх үзэгдэл. Иймд
\[
P=1-\frac{abc}{C_n^3}=1-\frac{6abc}{n(n-1)(n-2)}.
\]

\bod{2.4.}
$X_1, X_2$ нь харгалзан 1-р, 2-р буудагчийн онох тоо, $Q_i=1-P_i$ гэе. $X_i$ нь бином тархалттай:
$P(X_i=0)=Q_i^2$, $P(X_i=1)=2P_iQ_i$, $P(X_i=2)=P_i^2$. Буудагчид үл хамааран буудна.
\[
P(X_1>X_2)=P(X_1=1)P(X_2=0)+P(X_1=2)\bigl(P(X_2=0)+P(X_2=1)\bigr)
=2P_1Q_1Q_2^2+P_1^2(1-P_2^2).
\]

\bod{2.5.}
Цагаан бөмбөгний тоо хараасаа олон байх нь 8-аас 5, 6, 7 эсвэл 8 нь цагаан байх үзэгдэл.
\[
P=\frac{C_{15}^5C_{10}^3+C_{15}^6C_{10}^2+C_{15}^7C_{10}^1+C_{15}^8}{C_{25}^8}
=\frac{656370}{1081575}\approx 0.607.
\]

\bod{2.6.}
Бундан нь улаан өнгийн хөзөр тул $A\subset B$, иймд $C=A\cup B=B$.
$\ol{B}$ нь 4 мод бүгд хар байх үзэгдэл.
\[
P(C)=P(B)=1-\frac{C_{26}^4}{C_{52}^4}=1-\frac{14950}{270725}=\frac{787}{833}\approx 0.945.
\]

\bod{2.7.}
35 үсэг тус бүр нэг л картад бичигдсэн. “энхжин” үгэнд “н” үсэг хоёр удаа орж байгаа тул сугалсан 6 картаас энэ үгийг үүсгэх боломжгүй. Иймд аль ч тохиолдолд $P=0$.

(Хэрэв үгийн 6 үсэг бүгд өөр байсан бол эхний тохиолдолд $P=\frac{1}{A_{35}^6}$, байрыг сэлгэж болох үед $P=\frac{6!}{A_{35}^6}=\frac{1}{C_{35}^6}$ болно.)

\bod{2.8.}
Бөмбөг бүр $n$ хайрцгийн аль нэгэнд орох тул бүх боломжийн тоо $n^k$. Зэргэлдээ $k$ хайрцгийг $n-k+1$ янзаар сонгоно, тэдгээрт $k$ бөмбөгийг нэг нэгээр нь $k!$ янзаар байрлуулна. Иймд ($k\le n$ үед)
\[
P=\frac{(n-k+1)\,k!}{n^k}.
\]

\bod{2.9.}
Зэрэгцээ хоёр $P_1$: $1-(1-P_1)^2$. Зэрэгцээ хоёр $P_2$ ба тэдгээртэй зэрэгцээ $P_3$: $1-(1-P_2)^2(1-P_3)$. Эдгээр хоёр хэсэг ба $P_4$ цуваа холбогдсон тул
\[
P=\bigl[1-(1-P_1)^2\bigr]\bigl[1-(1-P_2)^2(1-P_3)\bigr]P_4 .
\]

\bod{2.10.}
Систем нь $n$ зэрэгцээ салаатай: нэг салаа нь ганц $P$, үлдсэн $n-1$ салаа нь цуваа холбогдсон $P_1$, $P$ (найдвар $P_1P$). Иймд
\[
R_n=1-(1-P)(1-P_1P)^{n-1}.
\]
$R_n\ge P_0$ нөхцөл нь $(1-P_1P)^{n-1}\le\dfrac{1-P_0}{1-P}$ болно. Логарифм авахад $\ln(1-P_1P)<0$ тул
\[
n\ge 1+\frac{\ln\dfrac{1-P_0}{1-P}}{\ln(1-P_1P)} .
\]
Энэ нөхцөлийг хангах хамгийн бага бүхэл $n$-ийг авна ($P\ge P_0$ бол $n=1$ хангалттай).

\bod{2.11.}
а) $n$ ширхэг зэрэгцээ хос цуваа холбогдсон: $R_a=\bigl(1-(1-P)^2\bigr)^n=P^n(2-P)^n$.

б) $n$ элементтэй хоёр цуваа салаа зэрэгцээ холбогдсон: $R_b=1-(1-P^n)^2=P^n(2-P^n)$.

$x=1-P$ гэвэл $(2-P)^n=(1+x)^n$, $2-P^n=2-(1-x)^n$ байна. Ньютоны биномоор задалбал
\[
(1+x)^n+(1-x)^n=2\bigl(1+C_n^2x^2+C_n^4x^4+\dots\bigr)\ge 2
\]
тул $(2-P)^n\ge 2-P^n$, өөрөөр хэлбэл $R_a\ge R_b$. Иймд а) систем илүү найдвартай ($n=1$ эсвэл $P\in\{0,1\}$ үед тэнцүү).

\bod{2.12.}
$A$, $A$, $B$ блокууд цуваа холбогдсон: $P_1=0.7\cdot0.7\cdot0.6=0.294$.
Схемд эхний $C$ блокийн зүүн үзүүр $-$ туйлтай холбогдсон тул 4 $C$ блок хүрээ үүсгэнэ. $K$ түлхүүр $C$ блокуудын хоорондох 3 цэгийн аль нэгэнд залгагдана. $p=0.8$.

1-р байрлал: гүйдэл нэг $C$-ээр эсвэл цуваа 3 $C$-ээр гарна,
\[
P_2=1-(1-p)(1-p^3)=1-0.2\cdot0.488=0.9024,\qquad P=1-0.294\cdot0.9024\approx0.735 .
\]
2-р байрлал: цуваа 2 $C$-тэй хоёр салаа зэрэгцээ холбогдсон,
\[
P_2=1-(1-p^2)^2=1-0.36^2=0.8704,\qquad P=1-0.294\cdot0.8704\approx0.744 .
\]
3-р байрлал нь 1-р байрлалтай адил: $P\approx0.735$.

\bod{2.13.}
Тэнцэх тохиолдлыг тооцохгүй гэж үзвэл өрөг бүрд аль нэг нь хожих магадлал $\frac12$. Бернулийн схемээр
\[
P(2:2)=C_4^2\Bigl(\frac12\Bigr)^4=\frac{6}{16}=0.375,\qquad
P(3:3)=C_6^3\Bigl(\frac12\Bigr)^6=\frac{20}{64}=0.3125 .
\]
Иймд 2:2 харьцаа илүү магадлалтай.

\bod{2.14.}
Нэг үйлдвэр дампуурахгүй байх магадлал $p=0.8$, $n=6$. $\nu$ нь дампуураагүй үйлдвэрийн тоо бол
\[
P_k=C_6^k\,0.8^k\,0.2^{6-k}.
\]
а) $P_2=15\cdot0.64\cdot0.0016=0.01536$.

б) $P(\nu>2)=1-P_0-P_1-P_2=1-0.000064-0.001536-0.01536=0.98304$.

\bod{2.15.}
$p=0.15$, $n=10$, $P_k=C_{10}^k\,0.15^k\,0.85^{10-k}$.

а) $P_3=120\cdot0.15^3\cdot0.85^7\approx0.130$.

б) $P(\nu\ge2)=1-P_0-P_1=1-0.85^{10}-10\cdot0.15\cdot0.85^9\approx1-0.197-0.347=0.456$.

\bod{2.16.}
Бином тархалтын хамгийн их магадлалтай тоо $k_0$ нь $np-q\le k_0\le np+p$ нөхцөлийг хангана. $p=0.95$, $q=0.05$, $k_0=100$ үед
\[
0.95n-0.05\le100\le0.95n+0.95
\;\Longrightarrow\;
104.26\le n\le105.32 .
\]
Иймд $n=105$ үр суулгах хэрэгтэй.

\bod{2.17.}
Зорчигч бүр 6 вагоны аль нэгийг сонгох тул бүх боломжийн тоо $6^{12}$.

а) $P=\dfrac{12!}{(2!)^6\,6^{12}}=\dfrac{7484400}{6^{12}}\approx0.0034$.

б) 6 вагонд 1, 2, 2, 0, 3, 4 хүн суух хуваарилалтыг $\frac{6!}{2!}=360$ янзаар сонгоно. Тогтсон хуваарилалтад зорчигчдыг $\frac{12!}{1!\,2!\,2!\,0!\,3!\,4!}=831600$ янзаар суулгана. Иймд
\[
P=\frac{360\cdot831600}{6^{12}}\approx0.138 .
\]

\bod{2.18.}
$A_i$ нь $i$-р шалгалтаа амжилттай өгөх үзэгдэл, $p_1=0.9$, $p_2=0.85$, $p_3=0.8$.
\[
P=p_1p_2p_3+p_1p_2q_3+p_1q_2p_3+q_1p_2p_3=0.612+0.153+0.108+0.068=0.941 .
\]

\bod{2.19.}
Буудалт бүрд онох магадлал $p_1=0.8$, $p_2=0.7$, $p_3=0.6$.

а) $P=1-q_1q_2q_3=1-0.2\cdot0.3\cdot0.4=0.976$.

б) $P=p_1p_2p_3+p_1p_2q_3+p_1q_2p_3+q_1p_2p_3=0.336+0.224+0.144+0.084=0.788$.

в) 3 дахь буудалтаар онох магадлал $p_3=0.6$.

г) 2-оос олонгүй алдах нь 3-уулаа алдахгүй байх үзэгдэл: $P=1-q_1q_2q_3=0.976$.

\bod{2.20.}
Англи, франц, немц хэлээр ярьдаг оюутнуудын олонлогийг $E$, $F$, $G$ гэе. Венийн диаграмаар:
зөвхөн англи 13, зөвхөн франц 7, зөвхөн немц 9, зөвхөн англи--франц 6, зөвхөн англи--немц 10, зөвхөн франц--немц 2, гурвуулаа 4. Иймд $|E\cup F\cup G|=51$, аль ч хэлээр ярьдаггүй 9 оюутан байна. $C_{60}^2=1770$.

а) Англиар ярьдаг, немцээр ярьдаггүй оюутан $13+6=19$: $P=\dfrac{C_{19}^2}{1770}=\dfrac{171}{1770}=\dfrac{57}{590}$.

б) Ядаж 2 хэлээр ярьдаг $6+10+2+4=22$: $P=\dfrac{22\cdot9}{1770}=\dfrac{33}{295}$.

в) Зөвхөн немц, францаар ярьдаг 2 оюутан: $P=\dfrac{1}{1770}$.

г) $P=1-\dfrac{C_{51}^2}{1770}=1-\dfrac{1275}{1770}=\dfrac{33}{118}$.

\bod{2.21.}
$P(A\cap B)=P(A)+P(B)-P(A\cup B)=\frac34+\frac35-\frac9{10}=\frac{9}{20}$.

а) $P(A/B)=\dfrac{9/20}{3/5}=\dfrac34$. \quad б) $P(B/A)=\dfrac{9/20}{3/4}=\dfrac35$.

$P(A/B)=P(A)$ тул $A$, $B$ үл хамаарах.

\bod{2.22.}
$P(A\cap B)=\frac25+\frac13-\frac35=\frac{2}{15}$.

а) $P(A/B)=\dfrac{2/15}{1/3}=\dfrac25$. \quad б) $P(B/A)=\dfrac{2/15}{2/5}=\dfrac13$. $P(A/B)=P(A)$ тул үл хамаарах.

в) $P(A\cap\ol{B})=P(A)-P(A\cap B)=\frac25-\frac2{15}=\frac4{15}$.

г) $P(\ol{A}\cap B)=P(B)-P(A\cap B)=\frac13-\frac2{15}=\frac15$.

д) $P(\ol{A}\cap\ol{B})=1-P(A\cup B)=1-\frac35=\frac25$.

\bod{2.23.}
$|\Omega|=6^3=216$. $B$: тэгш нүдтэй шоог 3 янзаар, тэгш нүдийг 3 янзаар, хоёр сондгой нүдийг $3^2$ янзаар сонгоно, $|B|=81$.
$A\cap B$: хоёр сондгой нүдийн ядаж нэг нь 5 байх ёстой, ийм хос $9-4=5$, $|A\cap B|=3\cdot3\cdot5=45$. $|A|=216-5^3=91$.

а) $P(A/B)=\dfrac{45}{81}=\dfrac59$. \quad б) $P(B/A)=\dfrac{45}{91}$.

\bod{2.24.}
$H_i$ нь бүтээгдэхүүн $i$-р шугамд үйлдвэрлэгдсэн байх, $A$ нь гологдол байх үзэгдэл. Гүйцэд магадлалын томьёогоор
\[
P(A)=0.45\cdot0.02+0.35\cdot0.04+0.20\cdot0.06=0.009+0.014+0.012=0.035 .
\]

\bod{2.25.}
II машины бүтээмжийг $x$ гэвэл I-ийнх $3x$, III-ынх $x/2$, нийт $4.5x$. Иймд $P(H_1)=\frac23$, $P(H_2)=\frac29$, $P(H_3)=\frac19$. Гүйцэд магадлалын томьёогоор
\[
P(A)=\frac23\cdot0.02+\frac29\cdot0.03+\frac19\cdot0.04=\frac{11}{450}\approx0.0244 .
\]

\bod{2.26.}
$H_1,H_2,H_3,H_4$ нь оюутан маш сайн, сайн, дунд, тааруухан бэлдсэн байх үзэгдлүүд, $A$ нь 3 асуултад бүгдэд нь хариулах үзэгдэл.
\[
P(H_i)=\tfrac{8}{20},\tfrac{6}{20},\tfrac{4}{20},\tfrac{2}{20};\qquad
P(A/H_i)=1,\ \frac{C_{35}^3}{C_{40}^3},\ \frac{C_{25}^3}{C_{40}^3},\ \frac{C_{10}^3}{C_{40}^3}
=1,\ \frac{6545}{9880},\ \frac{2300}{9880},\ \frac{120}{9880}.
\]
\[
P(A)=\frac{8\cdot9880+6\cdot6545+4\cdot2300+2\cdot120}{20\cdot9880}=\frac{127750}{197600}\approx0.6465 .
\]
Байесын томьёогоор

а) $P(H_2/A)=\dfrac{6\cdot6545}{127750}=\dfrac{39270}{127750}\approx0.307$.

б) $P(H_4/A)=\dfrac{2\cdot120}{127750}=\dfrac{240}{127750}\approx0.0019$.

\bod{2.27.}
$H_k$ нь өвчтөн $k$-р бүлгийн цустай байх үзэгдэл, $P(H_k)=0.337;\ 0.375;\ 0.209;\ 0.079$. $k$-р бүлгийн өвчтөнд $1,\dots,k$ бүлгийн цус таарах тул санамсаргүй донорын цус таарах магадлал
\[
P(A/H_1)=0.337,\quad P(A/H_2)=0.712,\quad P(A/H_3)=0.921,\quad P(A/H_4)=1 .
\]
а) Гүйцэд магадлалын томьёогоор
\[
P(A)=0.337\cdot0.337+0.375\cdot0.712+0.209\cdot0.921+0.079\cdot1\approx0.652 .
\]
б) 2 донор байхад ядаж нэгнийх нь цус таарах магадлал $1-(1-P(A/H_k))^2$ тул
\[
P=0.337\bigl(1-0.663^2\bigr)+0.375\bigl(1-0.288^2\bigr)+0.209\bigl(1-0.079^2\bigr)+0.079\approx0.819 .
\]

\bod{2.28.}
$A$ хайрцагт 5 сондгой, 4 тэгш, $B$ хайрцагт 3 сондгой, 2 тэгш хуудас байна. Хайрцаг бүрийг $\frac12$ магадлалтай сонгоно.

а) Эхний хуудас сондгой бол нөгөө хайрцгаас авна:
\[
P=\frac12\cdot\frac59\cdot\frac35+\frac12\cdot\frac35\cdot\frac59=\frac13 .
\]
б) Эхний хуудас тэгш бол мөн хайрцгаас буцаалтгүй авна:
\[
P=\frac12\cdot\frac49\cdot\frac38+\frac12\cdot\frac25\cdot\frac14=\frac1{12}+\frac1{20}=\frac{2}{15}.
\]
в) Байесын томьёогоор
\[
P=\frac{1/20}{2/15}=\frac38 .
\]

%======================================================================
\section*{II. Bertsekas, Tsitsiklis, Introduction to Probability, 2-р бүлэг}

\bod{2.2.}
Танаас бусад 499 зочны төрсөн өдөр тантай давхцах эсэх нь Бернулийн схем, $p=\frac1{365}$. Яг нэг зочин давхцах магадлал
\[
P=C_{499}^1\,\frac1{365}\Bigl(\frac{364}{365}\Bigr)^{498}\approx0.3487 .
\]
Пуассоны ойролцоо утга: $\lambda=np=\frac{499}{365}\approx1.367$, $P\approx\lambda e^{-\lambda}\approx0.3484$. Хоёр утга бараг тэнцүү байна.

\bod{2.3.}
Тоглолт бүрд Фишер хожих магадлал 0.4, Спасский хожих 0.3, тэнцэх 0.3.

(a) Фишер $l$-р тоглолтод хожихын тулд өмнөх $l-1$ тоглолт тэнцсэн байх ёстой ($l\le10$):
\[
P(\text{Фишер хожих})=\sum_{l=1}^{10}0.3^{\,l-1}\cdot0.4=0.4\cdot\frac{1-0.3^{10}}{0.7}\approx0.5714 .
\]
(b) $L$ нь тэмцээний үргэлжлэх тоглолтын тоо. $l<10$ үед эхний $l-1$ тэнцэж, $l$-д аль нэг нь хожно. $L=10$ нь эхний 9 тоглолт тэнцэх үзэгдэл:
\[
p_L(l)=\begin{cases}0.3^{\,l-1}\cdot0.7, & l=1,\dots,9,\\ 0.3^9, & l=10.\end{cases}
\]

\bod{2.14.}
$X$ нь $0,1,\dots,9$ утгыг $\frac1{10}$ магадлалтай авна.

(a) $Y=X \bmod 3$:
$P(Y=0)=P(X\in\{0,3,6,9\})=\frac4{10}$, $P(Y=1)=P(X\in\{1,4,7\})=\frac3{10}$, $P(Y=2)=P(X\in\{2,5,8\})=\frac3{10}$.

(b) $Z=5 \bmod (X+1)$: $X=0,1,2,3,4$ үед $Z=0,1,2,1,0$, $X\ge5$ үед $Z=5$. Иймд
\[
P(Z=0)=\tfrac2{10},\quad P(Z=1)=\tfrac2{10},\quad P(Z=2)=\tfrac1{10},\quad P(Z=5)=\tfrac5{10}.
\]

\bod{2.18.}
$X$ нь $2^a,2^{a+1},\dots,2^b$ утгуудыг $\frac1{b-a+1}$ магадлалтай авна.
\[
EX=\frac1{b-a+1}\sum_{k=a}^b2^k=\frac{2^{b+1}-2^a}{b-a+1},\qquad
EX^2=\frac1{b-a+1}\sum_{k=a}^b4^k=\frac{4^{b+1}-4^a}{3(b-a+1)},
\]
\[
DX=EX^2-(EX)^2=\frac{4^{b+1}-4^a}{3(b-a+1)}-\Bigl(\frac{2^{b+1}-2^a}{b-a+1}\Bigr)^2 .
\]

\bod{2.22.}
(a) Нэг орхилтод “нэг нь сүлд, нөгөө нь тоо” гарах магадлал $r=p(1-q)+q(1-p)$. Орхилтууд үл хамаарах тул орхилтын тоо $X$ нь $r$ параметртэй геометр тархалттай:
\[
P(X=k)=(1-r)^{k-1}r,\quad k=1,2,\dots;\qquad EX=\frac1r,\qquad DX=\frac{1-r}{r^2}=\frac{pq+(1-p)(1-q)}{\bigl(p(1-q)+q(1-p)\bigr)^2}.
\]
(b) Сүүлийн орхилтод эхний зоос сүлд гарах магадлал
\[
P(\text{СТ}/\{\text{СТ},\text{ТС}\})=\frac{p(1-q)}{p(1-q)+q(1-p)},
\]
$p=q$ үед энэ нь $\frac12$ байна.

\bod{2.25.}
(a) Нийт хариултын тоо $M=\sum_{k=1}^n m_k$ бөгөөд хариулт бүр ижил магадлалтай сонгогдоно:
\[
p_{I,J}(i,j)=\begin{cases}\dfrac1M, & j\le m_i,\\[4pt] 0, & \text{бусад үед}.\end{cases}
\]
Тухайн тархалтууд:
\[
p_I(i)=\sum_jp_{I,J}(i,j)=\frac{m_i}{M},\qquad p_J(j)=\sum_ip_{I,J}(i,j)=\frac{l_j}{M},
\]
энд $l_j$ нь $j$-р асуултад хариулсан ($m_i\ge j$) оюутны тоо.

(b) $i$-р оюутны $j$-р асуултын оноо $a$ (магадлал $p_{ij}$) эсвэл $b$ (магадлал $1-p_{ij}$) тул нийт онооны математик дундаж
\[
\sum_{j=1}^{m_i}\bigl(p_{ij}a+(1-p_{ij})b\bigr).
\]

\bod{2.32.}
$i$-р хосын эхний хүн амьд бол $X_i=1$, үгүй бол $X_i=0$; хоёр дахь хүнд мөн адил $Y_i$ гэе. Тэгвэл $S=\sum_{i=1}^mX_iY_i$ бөгөөд бүх хос ижил нөхцөлтэй тул
\[
E(S/A=a)=m\,E(X_1Y_1/A=a)=m\,P(X_1=1/A=a)\,P(Y_1=1/X_1=1,A=a).
\]
$A=a$ үед амьд $a$ хүн $2m$ хүнээс ижил магадлалтай сонгогдсон мэт тул
\[
P(X_1=1/A=a)=\frac{a}{2m},\qquad P(Y_1=1/X_1=1,A=a)=\frac{a-1}{2m-1}.
\]
Иймд
\[
E(S/A=a)=m\cdot\frac{a}{2m}\cdot\frac{a-1}{2m-1}=\frac{a(a-1)}{2(2m-1)} ,
\]
энэ нь $p$-ээс хамаарахгүй.

\bod{2.35.}
$E\,g(X,Y)=\sum_x\sum_yg(x,y)p_{X,Y}(x,y)$ ба $E(aX+bY)=aEX+bEY$ гэж батлах.

$Y=y$ үед $g(X,y)$ нь нэг санамсаргүй хэмжигдэхүүний функц тул
\[
E\bigl(g(X,Y)/Y=y\bigr)=\sum_xg(x,y)P(X=x/Y=y).
\]
Гүйцэд магадлалын томьёотой адилаар
\[
\begin{aligned}
E\,g(X,Y)&=\sum_yP(Y=y)\,E\bigl(g(X,Y)/Y=y\bigr)=\sum_y\sum_xg(x,y)P(Y=y)P(X=x/Y=y)\\
&=\sum_x\sum_yg(x,y)p_{X,Y}(x,y).
\end{aligned}
\]
$g(x,y)=ax+by$ үед
\[
\begin{aligned}
E(aX+bY)&=a\sum_xx\sum_yp_{X,Y}(x,y)+b\sum_yy\sum_xp_{X,Y}(x,y)\\
&=a\sum_xx\,p_X(x)+b\sum_yy\,p_Y(y)=aEX+bEY.
\end{aligned}
\]

\bod{2.39.}
$X_i$ нь $i$-р өдөр идэх өндөгний тоо, $X_i$ нь $1,\dots,6$ дээр жигд тархалттай, үл хамаарах.
\[
EX_i=\frac72,\qquad DX_i=\frac{6^2-1}{12}=\frac{35}{12}.
\]
$X=X_1+\dots+X_{10}$ тул
\[
EX=10\cdot\frac72=35,\qquad DX=10\cdot\frac{35}{12}\approx29.17 .
\]

\bod{2.44.}
$X$, $Y$ үл хамаарах бол $g(X)$, $h(Y)$ мөн үл хамаарах гэж батлах.

$U=g(X)$, $V=h(Y)$ гэе. $p_{X,Y}(x,y)=p_X(x)p_Y(y)$ тул
\[
P(U=u,V=v)=\sum_{\substack{g(x)=u\\h(y)=v}}p_X(x)p_Y(y)=\sum_{g(x)=u}p_X(x)\sum_{h(y)=v}p_Y(y)=P(U=u)P(V=v).
\]
Иймд $U$, $V$ үл хамаарах.

\bod{2.45.}
$X_1,\dots,X_n$ үл хамаарах, $X=X_1+\dots+X_n$, $EX=\mu$.

(a) $X_i$ нь $p_i$ параметртэй Бернулийн тархалттай бол $\sum p_i=\mu$ ба
\[
DX=\sum_{i=1}^np_i(1-p_i)=\mu-\sum_{i=1}^np_i^2 .
\]
$\sum p_i=\mu$ гэдгийг ашиглавал
\[
\sum_{i=1}^np_i^2=\sum_{i=1}^n\Bigl(p_i-\frac\mu n\Bigr)^2+\frac{\mu^2}{n},
\]
энэ нь $p_i=\frac\mu n$ үед хамгийн бага. Иймд $DX$ нь бүх $p_i=\frac\mu n$ үед хамгийн их.

(b) $X_i$ нь $p_i$ параметртэй геометр тархалттай бол $EX_i=\frac1{p_i}$, $DX_i=\frac{1-p_i}{p_i^2}$. $y_i=\frac1{p_i}$ гэвэл $\sum y_i=\mu$ ба
\[
DX=\sum_{i=1}^n\frac{1-p_i}{p_i^2}=\sum_{i=1}^ny_i(y_i-1)=\sum_{i=1}^ny_i^2-\mu .
\]
(a)-тай адилаар $\sum y_i^2$ нь $y_i=\frac\mu n$, өөрөөр хэлбэл $p_i=\frac n\mu$ үед хамгийн бага. Иймд $DX$ хамгийн бага байна.

\end{document}
